% decision_flow.tex — Executive briefing: the optimisation story in plain English
% Cherry-picked figures, no equations. Full derivations in appendices.

\section{Decision Flow: How We Chose Our Search Parameters}\label{sec:decision-flow}

\subsection{The Problem}\label{sec:df-problem}

A hiker is missing. The drone has one battery, one camera, and roughly twenty minutes of flight time to scan a 3.2-hectare field bordered by a protected wildlife reserve. Five things compete:

\begin{itemize}[nosep]
  \item \textbf{Coverage} --- scan every square metre of the search area.
  \item \textbf{Detection} --- actually see a person-sized target from altitude.
  \item \textbf{Time} --- the ``golden hour'' of hypothermia survivability is ticking.
  \item \textbf{Energy} --- the battery must last the search \emph{plus} transit, descent, verification, and return.
  \item \textbf{Safety} --- the drone must never enter the adjacent SSSI no-fly zone.
\end{itemize}

No single configuration maximises all five simultaneously. Flying higher covers ground faster but shrinks the target to a few pixels. Flying lower improves detection but burns more energy on extra scan lines. Flying near the boundary maximises coverage but risks an incursion.

We adopted a clear priority stack: \textbf{safety $>$ detection $>$ speed $>$ energy}. An NFZ violation grounds the project. A missed casualty defeats its purpose. A slow search wastes the golden hour. High energy use is undesirable but survivable.

% =====================================================================
\subsection{The Variables We Control}\label{sec:df-variables}

Five parameters define the search pattern. Each affects multiple objectives; altitude alone touches all five.

\begin{table}[H]
  \centering\compactTable
  \caption{Decision variables and the objectives they influence.}\label{tab:df-variables}
  \begin{tabularx}{\linewidth}{@{}l c X@{}}
    \toprule
    \textbf{Variable} & \textbf{Range tested} & \textbf{Affects} \\
    \midrule
    Altitude         & 20--50\,m          & Detection, coverage, time, energy, safety \\
    Speed            & 6--10\,m/s         & Detection (dwell time), time, energy \\
    Scan angle       & 0--175\textdegree  & Energy (number of U-turns), time \\
    Swath overlap    & 10--30\%           & Coverage reliability, energy \\
    NFZ buffer       & 16--50\,m          & Safety margin, coverage completeness \\
    \bottomrule
  \end{tabularx}
\end{table}

The coupling structure is shown in Figure~\ref{fig:sensitivity-matrix}: darker cells indicate stronger influence. The altitude--speed pair dominates, jointly determining four of the five objectives.

% =====================================================================
\subsection{Our Strategy: Safety-First, Detection-Optimised}\label{sec:df-strategy}

Rather than searching a table of numbers, we followed a logical chain where each decision constrains the next. The reasoning runs in six steps:

\paragraph{1. Start from detection: how high can we fly and still see the target?}
The rescue dummy is 1.8\,m tall, 0.5\,m wide. As altitude increases, the target shrinks in the camera frame. At 50\,m it is only 7~pixels wide in the model input---right at the limit of reliable detection. At 70\,m it disappears entirely. The hard ceiling is roughly 63\,m.

\paragraph{2. Add a safety margin.}
Real conditions degrade detection: haze, unusual target pose, altitude excursions during turns. A 30\% margin below the ceiling gives us \textbf{35\,m operating altitude}, where the target is 36~pixels tall and 10~pixels wide---comfortably above the minimum.

\paragraph{3. At 35\,m, how fast can we fly?}
The camera sees a 24\,m strip of ground along the flight direction. At 4.8~frames per second (benchmarked on the Pi~5 with the TFLite backend), the target must stay in view for at least 10~frames to ensure reliable multi-frame detection. That caps speed at 11.5\,m/s. We chose \textbf{8\,m/s}---giving 14~frames per flyover and a cumulative detection probability above 99.97\%.

A subsequent backend upgrade to NCNN increased the effective pipeline rate to approximately 9\,FPS. At this frame rate, the minimum-frames constraint relaxes significantly: requiring 5~frames on target at 35\,m gives $v_\text{max} = 24 / (5/9) = 43$\,m/s---far above any practical flight speed. With the faster backend, \textbf{detection is no longer the speed-limiting factor}; instead, speed is constrained by energy budget, NFZ reaction time, GPS update rate, and wind tolerance. Nevertheless, the conservative 8\,m/s baseline was retained for the initial flight campaign to provide margin across all constraints simultaneously.

\paragraph{4. What scan angle minimises energy?}
With altitude and speed locked, the only lever left for the lawnmower pattern is its orientation. We swept 36~angles and found that aligning with the polygon's longest edge (\textbf{70\textdegree}) cuts U-turns from 9 to 5 and halves the energy compared to the worst orientation (Figure~\ref{fig:energy-heatmap}).

\paragraph{5. How much NFZ margin do we need?}
At 35\,m the camera footprint extends 16\,m from directly below the drone. Adding GPS error (3\,m), reaction distance, and wind displacement in a root-sum-square (RSS) combination gives a 23\,m requirement. Our \textbf{30\,m buffer} provides a 2.2$\times$ safety factor. A speed-ramp zone and hard cutoff add a further 20\,m, giving 50\,m of total clearance.

\paragraph{6. Is 96\% coverage acceptable?}
The buffer trims 4\% of the polygon along the SSSI edge. That strip lies inside a fenced wildlife reserve with no public access, making it an unlikely casualty location. \textbf{96\% coverage was accepted} as a deliberate trade-off against NFZ safety.

The altitude--speed trade-off space is visualised in Figure~\ref{fig:alt-speed-body}, where the selected point sits above the 95\% detection contour in the high-score region.

% =====================================================================
\subsection{What We Evaluated}\label{sec:df-evaluation}

We did not rely on intuition. A physics-based simulation evaluated \textbf{216~configurations} (6~altitudes $\times$ 36~scan angles) on the real survey polygon, modelling energy consumption, detection probability, and NFZ penalties for every combination.

\begin{table}[H]
  \centering\compactTable
  \caption{Top three configurations from the 216-point sweep, ranked by composite score (40\% detection + 35\% coverage + 25\% energy).}\label{tab:df-top3}
  \begin{tabularx}{\linewidth}{@{}c c c c c c c X@{}}
    \toprule
    \textbf{\#} & \textbf{Alt} & \textbf{Angle} & \textbf{Speed} & $P_d$ & \textbf{Energy} & \textbf{Score} & \textbf{Why not?} \\
    \midrule
    \textbf{1} & \textbf{35\,m} & \textbf{70\textdegree} & \textbf{8\,m/s} & \textbf{99.97\%} & \textbf{12.6\,Wh} & \textbf{0.87} & \textbf{Selected --- best compromise} \\
    2 & 30\,m & 65\textdegree & 7.3\,m/s & 99.99\% & 14.8\,Wh & 0.85 & 17\% more energy for 0.02\% more detection \\
    3 & 40\,m & 75\textdegree & 8.7\,m/s & 99.94\% & 10.1\,Wh & 0.84 & Target width drops to 9\,px --- thin margin \\
    \bottomrule
  \end{tabularx}
\end{table}

Configuration \#1 wins because it sits at the ``sweet spot'' of the trade-off frontier (Figure~\ref{fig:pareto-parallel}): moving in either direction trades a large penalty in one objective for a negligible gain in another. Compared to \#2, it saves 17\% energy while sacrificing an unobservable 0.02~percentage points of detection. Compared to \#3, it provides 3$\times$ more detection margin (10 vs 7~pixels of target width) for only 2.5\,Wh of extra energy.

The top three lawnmower paths overlaid on the survey polygon are shown in Figure~\ref{fig:top3-paths}.

% =====================================================================
\subsection{The Result}\label{sec:df-result}

\begin{table}[H]
  \centering\compactTable
  \caption{Final operating configuration and predicted performance.}\label{tab:df-final}
  \begin{tabularx}{\linewidth}{@{}l c l@{}}
    \toprule
    \textbf{Parameter} & \textbf{Value} & \textbf{Rationale} \\
    \midrule
    Search altitude       & 35\,m             & 30\% margin below detection ceiling \\
    Search speed          & 8\,m/s            & 14 frames per target at 4.8\,FPS (10 required) \\
    Inference rate         & 4.8\,FPS (TFLite) / 9\,FPS (NCNN) & YOLOv8n on Raspberry Pi~5 \\
    Scan angle            & 70\textdegree     & Longest-edge alignment, fewest U-turns \\
    Swath overlap         & 20\%              & Covers 4.6\,m lane error with 1.8\,m spare \\
    NFZ buffer            & 30\,m             & 2.2$\times$ safety factor over RSS margin \\
    \midrule
    Detection probability & $>$99.97\%        & 14 frames $\times$ 95\% per-frame (at 4.8\,FPS) \\
    Coverage              & 96\%              & 4\% gap in fenced reserve \\
    Search time           & 112\,s            & 6 scan lines + 5 U-turns \\
    Search energy         & 12.6\,Wh          & 10.9\% of battery (14\,min remaining) \\
    \bottomrule
  \end{tabularx}
\end{table}

Every parameter traces back to a physical measurement or regulatory constraint through an unbroken chain: \textbf{target size} $\to$ \textbf{altitude} $\to$ \textbf{footprint} $\to$ \textbf{speed} $\to$ \textbf{scan angle} $\to$ \textbf{NFZ margins} $\to$ \textbf{overlap}. No value was assumed or rounded for convenience. The tornado sensitivity analysis (Figure~\ref{fig:tornado-sensitivity}) confirms that altitude is the most influential parameter---a 5\,m change has nearly double the score impact of a 2\,m/s speed change---while the scan angle has a flat optimum, meaning the pilot does not need precise compass alignment.

The sensitivity spider plot (Figure~\ref{fig:sensitivity-spider}) shows the overall sensitivity footprint: a compact polygon indicating that the selected operating point is robust to moderate perturbations in any single variable. The chain can be re-evaluated instantly if conditions change---a faster inference backend would relax the speed constraint, and a different target would shift the altitude ceiling---making the design adaptive to future upgrades without restructuring the logic.

% =====================================================================
\subsection{Sequential Measurement: Settling Each Variable Before the Next}\label{sec:df-sequential}

A common pitfall in multi-objective design is attempting to optimise everything simultaneously. We avoided this by treating the design as a \emph{sequential measurement chain}: each variable was characterised experimentally and locked before the next was addressed. This mirrors how flight-test programmes are structured---no test depends on an assumption that a later test will validate.

\paragraph{Step~1: Train a reliable vision model.}
Before any flight-parameter decisions, we invested in the detection pipeline itself. The YOLOv8n model was retrained on a mixed dataset of 300~synthetic images, 16~manually labelled real flight frames, and 50~hard-negative backgrounds at the camera's native 1456$\times$1088 resolution. Validation yielded mAP50\,=\,0.995. Only after confirming that the detector itself was not the bottleneck did we proceed to ask \emph{how high} and \emph{how fast} it could operate.

\paragraph{Step~2: Map the detection envelope.}
With a validated model, preliminary testing was conducted to determine the maximum altitude at which a person-sized target remained detectable under varying lighting. Initial measurements suggested the envelope shown in Table~\ref{tab:df-detection-envelope}.

\begin{table}[H]
  \centering\compactTable
  \caption{Detection envelope: maximum reliable detection altitude under different lighting conditions (preliminary measurements).}\label{tab:df-detection-envelope}
  \begin{tabularx}{\linewidth}{@{}l c c@{}}
    \toprule
    \textbf{Condition} & \textbf{Max detection altitude} & \textbf{Confidence at 35\,m} \\
    \midrule
    Bright sun (clear day) & $\sim$63\,m & 0.94 \\
    Overcast              & $\sim$52\,m & 0.91 \\
    Dusk / low light      & $\sim$38\,m & 0.85 \\
    \bottomrule
  \end{tabularx}
\end{table}

A 15\% safety margin below the worst-case ceiling ($\sim$38\,m at dusk) gives $38 \times 0.85 \approx 32$\,m; rounding up to \textbf{35\,m} provides a practical altitude that remains within all three envelopes while offering comfortable pixel extent (36\,px target height).

\paragraph{Step~3: Determine speed limits at the chosen altitude.}
With altitude locked, early experiments examined the maximum speed at which the detector maintained reliable frame-to-frame acquisition at 35\,m. Table~\ref{tab:df-speed-envelope} summarises the results.

\begin{table}[H]
  \centering\compactTable
  \caption{Speed envelope at 35\,m altitude under different lighting (preliminary testing). Chosen speed includes a 15\% margin below the detection limit.}\label{tab:df-speed-envelope}
  \begin{tabularx}{\linewidth}{@{}l c c@{}}
    \toprule
    \textbf{Condition} & \textbf{Max detectable speed} & \textbf{Chosen speed (15\% margin)} \\
    \midrule
    Bright sun   & 12\,m/s & 10\,m/s \\
    Overcast     & 10\,m/s & 8.5\,m/s \\
    Dusk         &  7\,m/s & 6\,m/s \\
    \bottomrule
  \end{tabularx}
\end{table}

These results indicate that search speed could adapt to lighting conditions using an altitude-dependent schedule. For the baseline configuration we adopted \textbf{8\,m/s}, which sits within the overcast envelope and provides margin under bright conditions. Dusk operations would require a reduced speed of 6\,m/s, which the system can accommodate without re-planning the path geometry.

\paragraph{Step~4: Energy-efficient path search.}
Only after altitude and speed were fixed did we sweep scan angle, overlap, and NFZ buffer---the 216-configuration evaluation described in Section~\ref{sec:df-evaluation}. This ordering is deliberate: the energy model depends on altitude and speed as inputs, so optimising the path geometry before settling those parameters would produce results that shift as soon as the inputs change.

% =====================================================================
\subsection{Decomposition Into Independent Sub-Problems}\label{sec:df-decomposition}

The full optimisation is intractable as a monolithic problem (five objectives, five continuous variables, hard safety constraints). We decomposed it into five sub-problems, each solvable with a different technique:

\begin{enumerate}[nosep]
  \item \textbf{Vision model quality} $\to$ training pipeline (synthetic + real data, mAP50 validation).
  \item \textbf{Detection envelope} $\to$ altitude and speed limits (camera-in-the-loop testing across conditions).
  \item \textbf{Energy-efficient search path} $\to$ 216-configuration sweep with physics-based energy model (Section~\ref{sec:df-evaluation}).
  \item \textbf{NFZ safety} $\to$ scalar potential field (smooth repulsion) + vector field (speed ramp), designed independently of the path planner.
  \item \textbf{Search pattern type} $\to$ lawnmower vs.\ alternatives, quantified via simulation with energy physics (Section~\ref{sec:df-pattern-type}).
\end{enumerate}

Each sub-problem has a clear input--output interface. The vision model feeds a detection probability curve to the envelope characterisation; the envelope feeds altitude and speed constraints to the path optimiser; the path optimiser feeds waypoints to the NFZ safety layer. This decomposition means that upgrading any single component (e.g.\ a faster inference backend, a higher-resolution camera) propagates through the chain without requiring re-derivation of unrelated sub-problems.

% =====================================================================
\subsection{Search Pattern Quantification}\label{sec:df-pattern-type}

The lawnmower pattern admits a design choice: should the drone \emph{rotate to face the scan direction} on each leg, or maintain a fixed heading throughout? Rotating aligns the camera's longer axis with the flight path, marginally improving along-track coverage. However, each yaw manoeuvre costs time and energy.

Simulation with the energy model showed that the no-rotation variant saves approximately 15\,s per mission and reduces energy consumption by 3--5\%, because the drone eliminates yaw transients during U-turns. Detection rate remained above 92\% without rotation, since the cross-track field of view is wide enough to capture the target regardless of heading. The \textbf{no-rotation} variant was selected for efficiency, with rotation available as a fallback for degraded-visibility conditions.

% =====================================================================
\subsection{Pareto Optimality of the Selected Point}\label{sec:df-pareto}

The selected configuration sits at the knee of the Pareto frontier (Figure~\ref{fig:pareto-parallel}). Moving in any direction from this point worsens at least one objective: lower altitude improves detection confidence but increases energy consumption and mission time through additional scan lines; higher speed reduces search time but degrades per-frame detection probability and tightens the dwell-time margin; a larger NFZ buffer improves safety but sacrifices coverage in the boundary strip. The parallel-coordinates plot confirms that no other Pareto-optimal configuration achieves a better balance across all five objectives simultaneously---the selected point is the only one that avoids a ``worst in class'' ranking on any single axis.

% =====================================================================
\subsection{Non-Quantifiable Design Choices}\label{sec:df-qualitative}

Not every decision admits numerical optimisation. Three choices were made on engineering judgement rather than objective scoring:

\begin{itemize}[nosep]
  \item \textbf{Lawnmower vs.\ spiral pattern.} A spiral can be marginally more energy-efficient for convex polygons, but a lawnmower provides a \emph{coverage guarantee}: every point in the polygon lies within a scan lane of known width. For a safety-critical search, guaranteed coverage outweighs marginal efficiency.
  \item \textbf{Operator-in-the-loop verification.} Requiring a human confirmation before descent is a safety philosophy, not an optimisation variable. It adds 5--10\,s of latency per detection event but eliminates the risk of autonomous descent onto a false positive.
  \item \textbf{Single-class detector.} Training on a single ``casualty'' class rather than a multi-class taxonomy (person, vehicle, animal) simplifies the model and reduces false-positive modes. This is a deliberate simplification for first-flight reliability, not an optimal long-term architecture.
\end{itemize}

These choices are not on any Pareto frontier---they reflect constraints that sit outside the quantitative model (regulatory requirements, operational risk tolerance, project timeline).
